% !PRINTED 603
\weberpage{603}
%<<CONT

By \S\,161, 4. we take in the class $(A_i)^{-1}$ an ideal $\mathfrak{m}$ coprime to $M$, and in

\weq{5}{\mathfrak{m}\,\mathfrak{a}_i = \overline{\alpha}}

we let $\mathfrak{a}_i$ run through the ideals of the class $\overline{A}_i$, and consequently $\overline{\alpha}$ through the ideals of the principal class divisible by $\mathfrak{m}$; if then

\[
N(\mathfrak{a}_i) \leqq \nu,
\]

then

\[
N(\mathfrak{m})\,N(\mathfrak{a}_i) = N(\alpha) \leqq \nu\,N(\mathfrak{m}).
\]

If, as before, $T$ denotes the number of ideals of the principal class divisible by $\mathfrak{m}$ whose norm is not greater than $t$, then, if we set $t = \nu N(\mathfrak{m})$,

\[
Z(\nu) = T,
\]

and from formula \S\,162, (6) there results

\weq{6}{Z(\nu) = g\nu + R_\nu\sqrt{\nu},}

in which $g$ denotes a number that is the same for all classes $\overline{A}_i$, and $R_\nu$ a function of $\nu$ that does not become infinite as $\nu$ increases to infinity.

By the meaning of $a_\nu$ we have

\weq{7}{a_\nu = Z(\nu) - Z(\nu - 1), \quad Z(0) = 0}

and consequently

\weq{8}{S_i = \sum_{0,\infty}^{\nu}\frac{Z(\nu) - Z(\nu - 1)}{\nu^s}.}

If we substitute the expression (6) for $Z(\nu)$, there follows:

\weq{9}{S_i = g\sum\frac{1}{\nu^s} + \sum R_\nu\sqrt{\nu}\left(\frac{1}{\nu^s} - \frac{1}{(\nu+1)^s}\right).}

The second term of this sum:

\[
\sum R_\nu\sqrt{\nu}\left(\frac{1}{\nu^s} - \frac{1}{(\nu+1)^s}\right)
\]

is unconditionally convergent so long as $s > \frac{1}{2}$, and is therefore a continuous function of $s$ which, for $s = 1$, becomes

\[
\sum_{1,\infty}\frac{R_\nu}{\sqrt{\nu}\,(\nu+1)}
\]

Indeed,

\[
\operatorname{Lim}_{\nu=\infty} \nu^{s+\frac{1}{2}} R_\nu \sqrt{\nu}\left(\frac{1}{\nu^s} - \frac{1}{(\nu+1)^s}\right) = s\operatorname{Lim} R_\nu
\]
%CONT>>
